\chapter{Touch and OpenHaptics Installation}
\label{app:openhaptics}

\begingroup
\small
\lstset{
  basicstyle=\ttfamily\footnotesize,
  aboveskip=0.4em,
  belowskip=0.4em
}

This appendix gives the complete device-specific installation used for the
3D Systems Touch. The validated stack is Ubuntu 22.04, ROS~2 Humble,
OpenHaptics 3.4, and Touch Driver 2025.12.10. Install the general local
environment from \filepath{README_ENV.md} first. Keep the Touch disconnected
until the driver installation is complete.

\section{Download and extract the vendor packages}

Download both Linux archives from the official 3D Systems support page:

\begin{center}
\url{https://support.3dsystems.com/s/article/OpenHaptics-for-Linux-Developer-Edition-v34?language=en_US}
\end{center}

\begin{itemize}
  \item \textbf{OpenHaptics Installer --- OpenHaptics for Linux v3.4};
  \item \textbf{Haptic Device Driver Installer --- Touch Device Driver
  v2025.12.10}.
\end{itemize}

Create a permanent working directory and extract both archives. The browser may
append a suffix such as \filepath{(1)} to a repeated download; adapt the archive
name when necessary.

\begin{lstlisting}[language=bash]
mkdir -p "$HOME/touch_driver"

tar -xzf "$HOME/Downloads/openhaptics_3.4-0-developer-edition-amd64.tar.gz" \
  -C "$HOME/touch_driver"
tar -xzf "$HOME/Downloads/TouchDriver_2025_12_10+1.tgz" \
  -C "$HOME/touch_driver"

ls "$HOME/touch_driver/TouchDriver_2025_12_10"
ls "$HOME/touch_driver/openhaptics_3.4-0-developer-edition-amd64"
\end{lstlisting}

Do not place these vendor archives inside the Git repository.

\section{Install dependencies and vendor software}

Install the build, terminal, and OpenGL dependencies required by the SDK and
its examples:

\begin{lstlisting}[language=bash]
sudo apt update
sudo apt install -y \
  build-essential freeglut3-dev libglu1-mesa-dev \
  libncurses5-dev mesa-utils
\end{lstlisting}

Install the Touch driver first. This copies the device libraries and installs
the USB permission rule:

\begin{lstlisting}[language=bash]
cd "$HOME/touch_driver/TouchDriver_2025_12_10"
sudo bash ./install_haptic_driver
sudo udevadm control --reload-rules
sudo udevadm trigger
\end{lstlisting}

Reconnect the Touch after this command. Reboot if requested. Then install the
OpenHaptics SDK:

\begin{lstlisting}[language=bash]
cd "$HOME/touch_driver/openhaptics_3.4-0-developer-edition-amd64"
sudo ./install
sudo ldconfig
\end{lstlisting}

The SDK installer places examples and documentation under
\filepath{/opt/OpenHaptics/Developer/3.4-0/}. The project also keeps the
extracted SDK root so that CMake can locate the exact headers and libraries.

\section{Configure the SDK path}

Define \filepath{OPENHAPTICS_ROOT} and verify the two files required by the
project build:

\begin{lstlisting}[language=bash]
export OPENHAPTICS_ROOT="$HOME/touch_driver/openhaptics_3.4-0-developer-edition-amd64"

test -f "$OPENHAPTICS_ROOT/usr/include/HD/hd.h" \
  && echo "OpenHaptics headers found"
test -f "$OPENHAPTICS_ROOT/usr/lib/libHD.so" \
  && echo "OpenHaptics HD library found"
\end{lstlisting}

Both messages must appear. Make the setting persistent for Bash terminals:

\begin{lstlisting}[language=bash]
echo 'export OPENHAPTICS_ROOT="$HOME/touch_driver/openhaptics_3.4-0-developer-edition-amd64"' \
  >> "$HOME/.bashrc"
source "$HOME/.bashrc"
\end{lstlisting}

For Zsh, add the same export line to \filepath{\textasciitilde/.zshrc} instead.

Some vendor executables may report missing legacy terminal libraries such as
\texttt{libtinfo.so.5} or \texttt{libncurses.so.5}. Install them only when that
error occurs:

\begin{lstlisting}[language=bash]
sudo apt install -y libtinfo5 libncurses5
\end{lstlisting}

\section{Validate the device before ROS}

Plug in the Touch. First compile and run the read-only vendor diagnostic:

\begin{lstlisting}[language=bash]
cd "$OPENHAPTICS_ROOT/opt/OpenHaptics/Developer/3.4-0/examples/HD/console/QueryDevice"
make
./QueryDevice
\end{lstlisting}

The program must identify the device, and the displayed pose and button values
must change when the stylus moves. If initialization fails, stop here and fix
the driver, USB connection, permissions, or missing library before using ROS.

Run the calibration example when the device requests calibration:

\begin{lstlisting}[language=bash]
cd "$OPENHAPTICS_ROOT/opt/OpenHaptics/Developer/3.4-0/examples/HD/console/Calibration"
make
./Calibration
\end{lstlisting}

Finally, an optional force-feedback test verifies bidirectional communication:

\begin{lstlisting}[language=bash]
cd "$OPENHAPTICS_ROOT/opt/OpenHaptics/Developer/3.4-0/examples/HD/console/FrictionlessSphere"
make
./FrictionlessSphere
\end{lstlisting}

Move the stylus gently; a virtual sphere should be perceptible. Keep a light
grip and stop immediately with \filepath{Ctrl+C} if the device behaves
unexpectedly.

\section{Build the project ROS~2 driver}

Activate Conda, ROS~2, and the SDK path before building the workspace:

\begin{lstlisting}[language=bash]
cd "$HOME/Stage_Lirmm/Diffusion-model-isaacsim"
conda activate microwave_dp
source /opt/ros/humble/setup.bash
export OPENHAPTICS_ROOT="$HOME/touch_driver/openhaptics_3.4-0-developer-edition-amd64"

cd ros2_WS
rosdep install --from-paths src --ignore-src -r -y
colcon build --symlink-install \
  --cmake-args -DOPENHAPTICS_ROOT="$OPENHAPTICS_ROOT"
source install/setup.bash
\end{lstlisting}

If the workspace was copied from another computer, old absolute paths may
remain in generated files. Remove only the generated directories and rebuild:

\begin{lstlisting}[language=bash]
cd "$HOME/Stage_Lirmm/Diffusion-model-isaacsim/ros2_WS"
rm -rf build install log
colcon build --symlink-install \
  --cmake-args -DOPENHAPTICS_ROOT="$OPENHAPTICS_ROOT"
source install/setup.bash
\end{lstlisting}

\section{Validate the ROS topics}

Start the Touch node in terminal 1:

\begin{lstlisting}[language=bash]
cd "$HOME/Stage_Lirmm/Diffusion-model-isaacsim"
conda activate microwave_dp
source /opt/ros/humble/setup.bash
source ros2_WS/install/setup.bash
ros2 launch touch_ros2_driver touch_rviz.launch.py
\end{lstlisting}

Inspect its outputs in terminal 2:

\begin{lstlisting}[language=bash]
source /opt/ros/humble/setup.bash
source "$HOME/Stage_Lirmm/Diffusion-model-isaacsim/ros2_WS/install/setup.bash"
ros2 topic hz /touch/pose
ros2 topic echo /touch/buttons
\end{lstlisting}

The pose must update continuously and both buttons must change the published
value. Stop both terminals with \filepath{Ctrl+C}. Do not continue to MuJoCo or
a physical robot if the topic freezes, the axes jump while the device is still,
or the node reports an OpenHaptics scheduler error.

\section{Daily use and common failures}

In every new terminal used with the Touch, activate the three software layers:

\begin{lstlisting}[language=bash]
cd "$HOME/Stage_Lirmm/Diffusion-model-isaacsim"
conda activate microwave_dp
source /opt/ros/humble/setup.bash
source ros2_WS/install/setup.bash
\end{lstlisting}

\begin{tabularx}{0.96\linewidth}{@{}>{\raggedright\arraybackslash}p{0.26\linewidth}>{\raggedright\arraybackslash}X@{}}
  \toprule
  \textbf{Symptom} & \textbf{Check} \\
  \midrule
  Device not detected & Reconnect USB, reload the udev rules, reboot, and run
  \filepath{QueryDevice} before ROS. \\
  Missing legacy terminal library & Install the two compatibility packages
  shown above. \\
  CMake cannot find \filepath{HD/hd.h} & Correct
  \filepath{OPENHAPTICS_ROOT}, then delete only \filepath{ros2_WS/build},
  \filepath{install}, and \filepath{log} before rebuilding. \\
  Scheduler or device-in-use error & Close other OpenHaptics examples and ROS
  nodes; only one process should own the device. \\
  ROS topics are absent & Source ROS~2 and \filepath{ros2_WS/install/setup.bash},
  then inspect the first terminal for a driver error. \\
  \bottomrule
\end{tabularx}

Only after the vendor diagnostics and ROS topics pass should the Touch be linked
to MuJoCo. Any later physical robot motion still requires the staged procedure
described in \cref{sec:physical-deployment}.

\endgroup
